% RESEARCH_DRAFT_SHA256: 0dfdd614d1cf42481a578b3777be0b0c408c9ecbea5758aaf17ee1f8d88b727e
% Layout and source-reference wrappers expanded for the web renderer.
\documentclass[11pt]{article}
\usepackage[margin=25mm]{geometry}
\usepackage{fontspec,unicode-math}
\setmainfont{Latin Modern Roman}
\setmathfont{Latin Modern Math}
\usepackage{amsmath,amssymb,mathtools,xurl,hyperref}
\hypersetup{colorlinks=true,urlcolor=blue,linkcolor=blue}
\setlength{\parindent}{0pt}
\setlength{\parskip}{5pt}
\setlength{\emergencystretch}{3em}
\newcommand{\sref}[2]{\hyperlink{source:#1:#2}{[S#2]}}
\newcommand{\cataloguescope}[1]{\paragraph{Source question.}#1}
\newcommand{\floor}[1]{\left\lfloor#1\right\rfloor}
\begin{document}

\section*{Frobenius roots: a terminating kernel certificate and an exact one-variable family}
\noindent catalogue entry; catalogue entry. Research attempt dated 11 October 2026.
\paragraph{Result in scope.}This manuscript contains a source-specific attempt and complete proofs of the stated auxiliary calculations. It makes no new solution or novelty claim for the entire target.
\section{Original problem and preserved definition}
The following definition and source list are reproduced from the repository. Any earlier inspection dates and reported literature claims in this preserved text describe that source record, not additional verification by this attempt.
% BEGIN PRESERVED DEFINITION
\subsection{Definitions and mathematical statement}

\paragraph{Shared notation.}Standard number systems are denoted by $\mathbb N,\mathbb Z,\mathbb Q,\mathbb R,\mathbb C$; index ranges, computational models and specialized notation follow the cited problem source and the definitions below.


In characteristic p, the Frobenius map sends $r$ to $r^p$; $I^{[p]}$ is generated by pth powers of elements of I. A unit F-module carries an isomorphism to its Frobenius pullback. A root is finite data whose iterated Frobenius maps generate the stated unit module, including local-cohomology modules in the regular-ring setting of the source.
\paragraph{Statement recorded in the supplied source.}
Local cohomology in characteristic p

If R is a regular ring of characteristic p, then the local cohomology modules have many nice finiteness properties; one can describe many of these properties via the theory of unit $F$-modules. Work on implementing these calculations in M2 was begun at the 2023 Minneapolis workshop.

Given an F-finite F-module, how can one efficiently compute a unit F-module? For example, given the map $\mathrm{Ext}^i(R/I,R)\to \mathrm{Ext}^i(R/I^{[p]},R)$, how can one find a root for the local cohomology module $H^i_I(R)$? \href{http://aimpl.org/macaulay2efie/2/}{[S2]}
\paragraph{Scope qualification.}
The regular-ring assumption and the supplied Ext map are retained. The source does not prescribe a ring encoding or a polynomial complexity bound.
\paragraph{Source status.}
The supplied catalogue labels this record partially solved. Its source-status metadata was checked against the supplied snapshot on 2026-10-11. This import does not independently establish current open status or mathematical correctness.
\subsection{Short English statement}
Can finite Frobenius-module data be converted efficiently into a root description of a local-cohomology module?
\subsection{Sources}
\begin{itemize}
\item[S1] Ulam AI and the UnsolvedMath Contributors, UnsolvedMath: A Curated Collection of Open Mathematics Problems, catalogue entry. Catalogue curation and metadata: CC BY 4.0. The supplied source question is retained; definitions and the English summary are editorial additions.
\url{https://huggingface.co/datasets/ulamai/UnsolvedMath}.
\item[S2] AIM Workshop Problem Lists, original problem source linked by the catalogue
\url{http://aimpl.org/macaulay2efie/2/}.
\end{itemize}
% END PRESERVED DEFINITION
\subsection{Verbatim supplied catalogue question}
Local cohomology in characteristic p

If R is a regular ring of characteristic p, then the local cohomology modules have many nice finiteness properties; one can describe many of these properties via the theory of unit $F$-modules. Work on implementing these calculations in M2 was begun at the 2023 Minneapolis workshop.

Given an F-finite F-module, how can one efficiently compute a unit F-module? For example, given the map $\mathrm{Ext}^i(R/I,R)\to \mathrm{Ext}^i(R/I^{[p]},R)$, how can one find a root for the local cohomology module $H^i_I(R)$?

\section{Attempt and results}
We first state the hypotheses needed for a kernel algorithm, then test it on
an explicit regular ring. These standard constructions overlap with the
supplied research summary; complete proofs below replace reliance on that
summary. They do not supply the requested efficient general implementation.

\paragraph{A fixed-module kernel iteration under exact Frobenius pullback.}
Let $R$ be Noetherian, $F$ an exact additive functor on its module category,
and $\beta:M\to F(M)$ a map with $M$ finitely generated. In the application,
$F$ is Frobenius pullback. Write
$\beta_e=F^{e-1}(\beta)\cdots F(\beta)\beta$ for $e\ge1$ and
$\beta_0=1_M$. Put $K_e=\ker\beta_e$. Since
$\beta_{e+1}=F(\beta_e)\beta$, exactness gives
\[
K_0=0,\qquad K_{e+1}=\beta^{-1}(F(K_e)).
\]
Here $F(K_e)$ is identified with its injected image in $F(M)$.
The kernels form an increasing chain because each iterate factors through
the preceding one. Noetherianity gives $K_t=K_{t+1}$ for some finite $t$.
The displayed recursion then implies by induction that all later kernels
equal $K_t=K$. Further, $\beta(K)\subseteq F(K)$, and exactness identifies
$F(M/K)$ with $F(M)/F(K)$. Thus the induced map
$\bar\beta:M/K\to F(M/K)$ has kernel $K_{t+1}/K_t=0$.

The original and quotient directed systems have the same direct limit.
Indeed their stage maps are the quotient maps $F^j(M)\to F^j(M/K)$, so they
induce a surjection on limits. If a stage-$j$ element maps to zero in the
quotient limit, at some later stage $l$ its image lies in $F^l(K)$.
The map from stage $l$ to stage $l+t$ is $F^l(\beta_t)$ and annihilates
$F^l(K)$ by exactness, so the element is already zero in the original limit.
This proves injectivity. It also proves that an equality of successive
computed kernels is a valid stopping certificate. ACC alone supplies no
numerical bound on $t$ or on the sizes of the presentations.

\paragraph{Exactness in a stated regular example.}
For $R=\mathbb F_p[x]$, every polynomial has a unique expression
$\sum_{i=0}^{p-1}x^i f_i(x)^p$, by sorting exponents modulo $p$ and using that
coefficients have unique $p$th roots in $\mathbb F_p$.
Thus $R$, regarded as a right $R$-module by $a\cdot r=ar^p$, is free with
basis $1,x,\ldots,x^{p-1}$. Tensoring with this module is exact: after choosing
the basis it is a direct sum of $p$ copies of the underlying exact sequence.
This establishes the needed hypothesis here without silently using flatness
for arbitrary rings. It also identifies $F(R/(x^m))=R/(x^{pm})$.

\paragraph{A complete scalar test family.}
Fix $m\ge1$ and define $\beta:R/(x^m)\to R/(x^{pm})$ by multiplication by
$x^a$, where $a\ge(p-1)m$ is an integer. This inequality makes the map
well defined. At stage $e$ the module is $R/(x^{mp^e})$, and the transition
multiplies by $x^{ap^e}$. The $e$-fold map from stage zero multiplies by
$x^{D_e}$, where $D_e=a(p^e-1)/(p-1)$. Consequently a monomial $x^j$,
$0\le j<m$, belongs to $K_e$ exactly when $j+D_e\ge mp^e$.
If $a=(p-1)m$, then $D_e=m(p^e-1)$ and no such $j$ is killed. The module
at stage $e$ embeds in $R_x/R$ by
\[
\bar v\longmapsto v/x^{mp^e}+R.
\]
The map is well defined and injective, since its kernel consists exactly of
multiples of $x^{mp^e}$, and it respects the transition maps. Every class of
a fraction with power-of-$x$ denominator occurs once $mp^e$ exceeds that
power, so the direct limit is $R_x/R$.
If $a>(p-1)m$, then
$D_e-mp^e=(a/(p-1)-m)p^e-a/(p-1)$ eventually becomes nonnegative.
All of $M$ is then in $K_e$, and the stable quotient and direct limit are zero.
Thus a single extra exponent can change a nonzero injective root to an
ultimately killed generating module. An implementation that simply calls
all finite generating data a root would miss this distinction.

\section{Context, adversarial audit, and gap}
The source names a particular Ext map and a local-cohomology module. No proof
identifying that map's direct limit with the stated local cohomology, or
algorithm constructing its presentations, is supplied here; such identification
must not be inferred from the abstract kernel lemma alone. The original AIMPL
page failed retrieval, and the primary workshop report supplies only context.
The audit checked functor exactness, the direction of the quotient map,
stage indices, equality as a stopping condition, and the threshold exponent.
No numerical experiments, package implementation, or complexity bound is claimed.

\section{Process and contributions}
Mathematics and model-assisted adversarial self-audit: OpenAI gpt-6-astra, ultra reasoning, as dispatched for this research batch; exact serving revision not exposed. Codex (GPT-6; exact serving revision not exposed) materially assisted with orchestration and format review. No human mathematical authorship, mathematical edits, or independent human proof checking is asserted.
Contribution prompt A was applied to the full definition and supplied source row, then prompt B was applied adversarially to the derived results. The item-specific audit above records the scope checks and rejected extrapolations. Local assembly and structural checks do not establish mathematical correctness. The source catalogue status is preserved in TOP\_PROBLEM and is separate from this attempt status.
\paragraph{Integration review.} Codex (GPT-6; exact serving revision not exposed) reviewed the written arguments and expanded layout and source-reference wrappers. Repository metadata and the production web renderer were checked using Python 3.13 and JavaScript with MathJax 3.2.2. This model-assisted review is not independent human proof verification; no PDF compilation or formal proof checking is asserted.

\section{Completion Estimate}
\noindent\textbf{COMPLETION: 0\%}
The proved material is a rederivation of known or auxiliary facts rather than a substantiated new advance on the original target. The kernel theorem is conditional on effective exact pullback; efficient general root extraction and the canonical Ext-to-local-cohomology identification are not established here.
\section{References}
The complete original Sources subsection above is retained, including catalogue attribution and licence notices. The additional list records primary sources actually consulted for this attempt on 11 October 2026; uncited source links reproduced above were not all rechecked.
\begin{itemize}
\item AIM, Macaulay2: expanded functionality and improved efficiency, official workshop summary; local-cohomology, DG, and direct-summand discussion read on 11 October 2026. The original AIMPL problem pages failed retrieval; their wording is available here only from the supplied snapshot. \url{https://aimath.org/pastworkshops/macaulay2efierep.pdf}.
\end{itemize}
\section{Final status}
\textbf{Unresolved in this attempt.} The kernel theorem is conditional on effective exact pullback; efficient general root extraction and the canonical Ext-to-local-cohomology identification are not established here.
\end{document}
